\documentclass[a4paper, serif, 10pt]{moderncv}

%% moderncv
% style and color
\moderncvstyle{banking}
\moderncvcolor{black}

%% page layout/margins
\usepackage[top=2.5cm, bottom=2.5cm, left=1.5cm, right=1.5cm]{geometry}

\nopagenumbers{}

%% Babel config for Spanish language
% \usepackage[spanish]{babel} has some conflicting issues with moderncv
% so we have to use enable the following options to make it work
%
% ref:
% - https://tex.stackexchange.com/a/140161/36007
\usepackage[spanish,es-lcroman]{babel}

%% fontspec
\usepackage{fontspec}

\IfFontExistsTF{Linux Libertine}{
  \setmainfont[Ligatures={TeX, Common}, Numbers=OldStyle]{Linux Libertine}
}{}
\IfFontExistsTF{Linux Libertine O}{
  \setmainfont[Ligatures={TeX, Common}, Numbers=OldStyle]{Linux Libertine O}
}{}

%% CTeX
% CJK support, used to show CJK characters in the resume
%
% - fontset=none: disable builtin fontset but instead set the CJK font manually
% - heading=false: disable ctex heading
% - punct=kaiming: use kaiming punctuations styles for CJK
% - scheme=plain: use plain scheme, do not override \normalsize font size
% - space=auto: space settings for CJK characters
%
% ref:
% - http://ctan.mirrorcatalogs.com/language/chinese/ctex/ctex.pdf
\usepackage[UTF8, heading=false, punct=kaiming, scheme=plain, space=auto]{ctex}

\IfFontExistsTF{Noto Serif CJK SC}{
  \setCJKmainfont{Noto Serif CJK SC}
}{}
\IfFontExistsTF{Noto Sans CJK SC}{
  \setCJKsansfont{Noto Sans CJK SC}
}{}

%% line spacing
\usepackage{setspace}
\setstretch{1.125}

%% URL styling - use normal text instead of monospace
\urlstyle{same}

% Auto-underline all links
% Defer until \begin{document} so hyperref is already loaded
\AtBeginDocument{%
  \let\oldhref\href
  \renewcommand{\href}[2]{\oldhref{#1}{\underline{#2}}}%
}

%% Patch \httplink and \httpslink to handle full URLs
% The original moderncv macros always prepend http:// or https:// to URLs.
% This causes issues when the URL already has a protocol (e.g., https://example.com)
% resulting in malformed URLs like https://https://example.com.
% This patch checks if the URL already contains "://" and uses it directly if so.
% Note: We use \str_set:Nx (x-expansion) to expand macros like \@homepage first.
\ExplSyntaxOn
\str_new:N \g_yamlresume_url_str

\RenewDocumentCommand{\httpslink}{O{}m}{
  \str_set:Nx \g_yamlresume_url_str {#2}
  \str_if_in:NnTF \g_yamlresume_url_str {://}
    { \url{#2} }
    { \url{https://#2} }
}

\RenewDocumentCommand{\httplink}{O{}m}{
  \str_set:Nx \g_yamlresume_url_str {#2}
  \str_if_in:NnTF \g_yamlresume_url_str {://}
    { \url{#2} }
    { \url{http://#2} }
}
\ExplSyntaxOff

\name{Laura Gómez Fernández}{}
\title{Product Manager Senior | Estrategia de producto y transformación digital}
\phone[mobile]{+34 612 345 678}
\email{laura@lauragomez.es}
\homepage{https://lauragomez.es}

\address{Calle de Alcalá 45, 3º D, Madrid, Comunidad de Madrid, España, 28001}{}{}

\extrainfo{{\small \faLinkedin}\ \href{https://www.linkedin.com/in/lauragomezpm}{@lauragomezpm} {} {} {} • {} {} {}
{\small \faMedium}\ \href{https://medium.com/@lauragomezpm}{@lauragomezpm}}

\begin{document}

\maketitle

\section{Información básica}

\cvline{}{Product Manager Senior con más de 6 años de experiencia en el ciclo de vida completo de productos digitales, desde la investigación de usuario hasta el lanzamiento y crecimiento.
Especializada en estrategia de producto, definición de hoja de ruta (roadmap), OKRs y descubrimiento de producto en entornos ágiles.
Experiencia en sectores fintech, movilidad y ecommerce, liderando equipos multifuncionales y colaborando estrechamente con diseño, ingeniería y negocio.
Orientada a datos, con toma de decisiones basada en métricas, analítica de producto, experimentación y A/B testing.}
\section{Educación}

\cventry{sept 2020–jul 2021}
        {Maestría, Dirección de Producto, Puntuación: 8.9}
        {IE Business School}
        {\url{https://www.ie.edu}}
        {}
        {Programa intensivo centrado en el desarrollo de habilidades de liderazgo de producto y toma de decisiones basada en datos.
Proyecto final: diseño de una solución de banca digital para jóvenes, con validación de usuario y prototipo funcional.
\textbf{Cursos}: Estrategia de Producto, Analítica Digital, Design Thinking, Liderazgo Ágil, Finanzas para Product Managers}

\cventry{sept 2013–jul 2017}
        {Licenciatura, Administración y Dirección de Empresas, Puntuación: 8.2}
        {Universidad Autónoma de Madrid}
        {\url{https://www.uam.es}}
        {}
        {Formación en gestión empresarial con especial interés en marketing digital y estrategia.
Trabajo de fin de grado sobre la adopción de plataformas digitales en el sector retail.
\textbf{Cursos}: Dirección Estratégica, Marketing Estratégico, Estadística Aplicada, Finanzas Corporativas, Comportamiento del Consumidor}
\section{Experiencia laboral}

\cventry{jun 2023 hasta la fecha}
        {Senior Product Manager}
        {Mobiliti}
        {\url{https://mobiliti.io}}
        {}
        {Lidero la estrategia de producto de la app de movilidad compartida con más de 2M de usuarios en España y América Latina.
Defino la hoja de ruta del producto alineada con los OKRs de la compañía, priorizando iniciativas de retención y monetización.
Coordino tres squads multifuncionales de diseño, ingeniería y datos para lanzar funcionalidades de pago, geolocalización y fidelización.
Impulso la experimentación continua mediante A/B testing y análisis de métricas de producto para reducir el churn y mejorar el LTV.
Colaboro con marketing y atención al cliente para incorporar insights de usuario y elevar la satisfacción (CSAT/NPS).
\textbf{Palabras clave}: Roadmap, OKRs, Metodologías ágiles, Analítica de producto, A/B testing, Retención}

\cventry{ene 2021–may 2023}
        {Product Manager}
        {Finpagos}
        {\url{https://finpagos.com}}
        {}
        {Gestioné el producto de pasarela de pagos B2B, desde la identificación de oportunidades hasta el lanzamiento y adopción.
Rediseñé el flujo de onboarding de comercios, logrando reducir el tiempo de alta en un 40\% y aumentando la conversión.
Implementé un proceso de descubrimiento de producto con entrevistas, tests de usabilidad y prototipado rápido.
Colaboré con ingeniería en la priorización del backlog y la definición de criterios de aceptación.
Presenté métricas y resultados a stakeholders, apoyándome en dashboards con Tableau y Amplitude.
\textbf{Palabras clave}: Fintech, Pasarela de pagos, Descubrimiento de producto, Backlog, Tableau, Amplitude}

\cventry{mar 2018–dic 2020}
        {Product Owner}
        {RetailWise}
        {\url{https://retailwise.io}}
        {}
        {Actué como product owner de una plataforma de comercio electrónico para clientes de retail y gran consumo.
Definí la visión y el backlog del producto en colaboración con stakeholders de negocio y tecnología.
Coordiné los sprints de desarrollo y las entregas a producción en ciclos quincenales con metodología Scrum.
Implementé mejoras en el proceso de checkout que incrementaron la tasa de conversión media en un 15\%.
Realicé análisis de competencia y estudios de mercado para fundamentar las decisiones de producto.
\textbf{Palabras clave}: Ecommerce, Scrum, Product Owner, Conversión, Investigación de mercado}
\section{Idiomas}

\cvline{Español}{Competencia nativa o bilingüe \hfill \textbf{Palabras clave}: Lengua materna, Comunicación, Redacción}
\cvline{Inglés}{Competencia profesional plena \hfill \textbf{Palabras clave}: C1 Advanced, Inglés de negocio, Reuniones internacionales}
\cvline{Catalán}{Competencia limitada de trabajo \hfill \textbf{Palabras clave}: Comprensión lectora, Conversación básica}
\section{Competencias}

\cvline{Gestión de Producto}{Experto \hfill \textbf{Palabras clave}: Roadmaps, OKRs, Priorización, Story Mapping, Product Discovery, Go-to-Market}
\cvline{Metodologías Ágiles}{Experto \hfill \textbf{Palabras clave}: Scrum, Kanban, Lean Startup, Design Sprint}
\cvline{Análisis de Datos}{Avanzado \hfill \textbf{Palabras clave}: SQL, Amplitude, Mixpanel, Google Analytics, A/B testing}
\cvline{UX / UI}{Intermedio \hfill \textbf{Palabras clave}: Figma, Prototipado, User Research, Usabilidad}
\cvline{Herramientas de Producto}{Experto \hfill \textbf{Palabras clave}: Jira, Confluence, Notion, Trello}
\section{Premios}

\cventry{nov 2022}
        {Finpagos}
        {Premio al Mejor Lanzamiento de Producto}
        {}
        {}
        {Reconocimiento interno al equipo de producto por el lanzamiento de la nueva pasarela de pagos, que alcanzó un crecimiento del 32\% en transacciones en el primer trimestre.}

\cventry{feb 2024}
        {Mobiliti}
        {Reconocimiento al Liderazgo en Producto}
        {}
        {}
        {Destacada por mi contribución a la mejora de la retención de usuarios y la consolidación de la estrategia de producto de la compañía.}
\section{Certificados}

\cventry{jun 2022}
        {Scrum.org}
        {Professional Scrum Product Owner (PSPO I)}
        {\url{https://www.scrum.org/certifications/pspo-i}}
        {}
        {}

\cventry{mar 2020}
        {Product School}
        {Product Management Certification}
        {\url{https://productschool.com}}
        {}
        {}

\cventry{abr 2021}
        {Google}
        {Google Analytics Individual Qualification}
        {\url{https://skillshop.exceedlms.com}}
        {}
        {}
\section{Publicaciones}

\cventry{nov 2022}
        {Cómo descubrir oportunidades de producto a través de la experimentación}
        {Medium}
        {\url{https://medium.com/@lauragomezpm/como-descubrir-oportunidades-de-producto}}
        {}
        {Artículo en el que comparto un marco práctico para priorizar experimentos de producto basados en datos y aprendizaje continuo.}

\cventry{jul 2023}
        {Claves para construir un roadmap de producto orientado a resultados}
        {LinkedIn Articles}
        {\url{https://www.linkedin.com/pulse/claves-roadmap-producto-laura-gomez}}
        {}
        {Guía práctica para conectar la visión de producto con la ejecución táctica del equipo y la estrategia de negocio.}
\section{Referencias}

\cventry{\emaillink[marta.lopez@mobiliti.io]{marta.lopez@mobiliti.io}}
        {Directora de Producto en Mobiliti}
        {Marta López Sánchez}
        {+34 699 123 456}
        {}
        {Laura es una profesional excepcional, con una gran capacidad para priorizar, comunicar y ejecutar estrategias de producto. Su liderazgo ha sido clave para el crecimiento de nuestro producto.}
\section{Proyectos}

\cventry{sept 2022 hasta la fecha}
        {Comunidad online para mentorear a product managers en activo.}
        {MentorPM}
        {\url{https://mentorpm.es}}
        {}
        {Creé una comunidad de producto y crecimiento con más de 500 miembros activos.
Organizo mentorías mensuales con product managers senior de empresas tecnológicas.
Desarrollo una newsletter semanal con aprendizajes de producto y liderazgo.
\textbf{Palabras clave}: Comunidad, Mentoría, Liderazgo, Newsletter}

\cventry{ene 2023–may 2023}
        {Prototipo en Figma de una aplicación móvil de bienestar y hábitos diarios.}
        {App de hábitos saludables}
        {\url{https://github.com/lauragomez/habitos}}
        {}
        {Diseñé un MVP en Figma con pruebas de usuario para validar la propuesta de valor.
Definición de métricas de éxito y test de usabilidad con 15 participantes.
\textbf{Palabras clave}: UX, Figma, Prototipado}
\section{Intereses}

\cvline{Lectura}{Producto, Marketing, Novela histórica}
\cvline{Running}{Maratones, Trail}
\cvline{Tecnología}{Inteligencia artificial, Startups, Low-code}
\section{Voluntariado}

\cventry{mar 2022 hasta la fecha}
        {Mentora de producto}
        {Mujeres Tech España}
        {\url{https://mujerestech.es}}
        {}
        {Acompaño a mujeres que quieren iniciarse en gestión de producto mediante sesiones de mentoría mensuales.
Colaboro en el diseño de talleres de desarrollo profesional y de liderazgo.}

\end{document}